\documentclass[10pt,a4paper]{article}
\usepackage{amsmath,amssymb,amsthm,graphicx}
\usepackage[a4paper,margin=28mm]{geometry}
\usepackage{amsmath,amssymb,amsthm,aliascnt,mathtools}
\usepackage{hyperref}
\hypersetup{colorlinks=true,linkcolor=cyan,citecolor=red,urlcolor=cyan,filecolor=cyan,
 pdftitle={A rank bound for bases and circuits in binary matroids},
 pdfauthor={},pdfkeywords={Binary matroid; bases; circuits}}
\numberwithin{equation}{section}
\newtheorem{theorem}{Theorem}[section]
\newaliascnt{lemma}{theorem}
\newtheorem{lemma}[lemma]{Lemma}
\aliascntresetthe{lemma}
\providecommand*{\lemmaautorefname}{Lemma}
\newaliascnt{proposition}{theorem}
\newtheorem{proposition}[proposition]{Proposition}
\aliascntresetthe{proposition}
\providecommand*{\propositionautorefname}{Proposition}
\newaliascnt{corollary}{theorem}
\newtheorem{corollary}[corollary]{Corollary}
\aliascntresetthe{corollary}
\providecommand*{\corollaryautorefname}{Corollary}
\renewcommand*{\theoremautorefname}{Theorem}
\renewcommand*{\sectionautorefname}{Section}
\renewcommand*{\subsectionautorefname}{Section}
\DeclareMathOperator{\GF}{GF}
\DeclareMathOperator{\cl}{cl}
\DeclareMathOperator{\spanof}{span}
\newcommand{\BB}{\mathcal B}
\newcommand{\CC}{\mathcal C}
\newcommand{\FF}{\GF(2)}
\newcommand{\wloc}[1]{w_{#1}(M;B)}
\newcommand{\dloc}[1]{d_{#1}(M;B)}
\newcommand{\ind}{\mathbf 1}
\title{Joint near-spanning circuit counts through an exact folded-cube edge correspondence}
\author{Houshan Fu}
\date{arxiv:2610.11829v1; CC BY 4.0}
\begin{document}
\maketitle
\noindent\textit{Editorial note.} Complete original author body and embedded bibliography retained; exactly one selected primary per task. No mathematical alteration, added assumption, proof repair or correctness certification.\par\medskip
\noindent\textit{Allowed standard prerequisites (editorial).}\begin{itemize}
\item Binary matroids, fundamental circuits and cube edge isoperimetry. All native matroid conventions, weighted circuit definitions, contraction identities, local bounds and original bibliography retained.. The selected comparison or exact joint count is retained in the full counted author argument.
\end{itemize}\medskip
\section{Introduction}

Comparing the numbers of maximal independent sets and minimal dependent
sets is a natural enumerative problem in matroid theory. For the cycle
matroid of a connected graph, this is precisely the comparison between
spanning trees and cycles.
For a matroid \(M\), let \(b(M)\), \(d(M)\), and \(r(M)\) denote
the number of bases, the number of circuits, and the rank, respectively.
Welsh \cite[p.~288, Exercise~5]{Welsh1976} posed an average version
of the comparison between bases and circuits for matroids of fixed
rank and size. For graphs, Quirk and Seymour proved that a connected
simple graph has at least as many spanning trees as cycles; Welsh
records their result as a private communication \cite[p.~287]{Welsh1976}.

Oxley \cite[Theorem~2.1]{Oxley1983} proved that every simple binary
matroid satisfies
\begin{equation}\label{eq:oxley}
 b(M)\ge2d(M),
\end{equation}
with equality precisely when deleting all coloops leaves the Fano
matroid \(F_7\). For nonempty members of this class with no coloops,
he also proved \(b(M)>6(r(M)+1)d(M)/19\)
\cite[Theorem~3.1]{Oxley1983}. The referee of that paper conjectured
that \(2b(M)\ge(r(M)+1)d(M)\).
The conjecture is recorded as Conjecture~15.5.16 in \cite{Oxley2011}.
Our first theorem proves it and determines the equality case.

\begin{theorem}\label{thm:rank}
Let \(M\) be a nonempty simple binary matroid with no coloops. Then
\[
 2b(M)\ge(r(M)+1)d(M).
\]
Equality holds if and only if \(M\cong F_7\).
\end{theorem}

For ranks at least four, this strengthens the rank-independent bound
\eqref{eq:oxley}. The coefficient is sharp for the class as a whole,
since \(b(F_7)=28\) and \(d(F_7)=14\).
The exclusion of coloops is necessary: adjoining coloops increases
the rank without changing either count. For example, let \(M\) be
\(F_7\) with one added coloop. Then \(r(M)=4\) and
\(2b(M)=56<70=(r(M)+1)d(M)\).

Related circuit-counting results include lower bounds for connected
matroids due to Oxley and Wang \cite[Lemmas~13 and~14]{OxleyWang2019}
and for circuits containing a fixed element in a \(3\)-connected matroid
with at least four elements due to Cho \cite[Theorem~3.4.1]{Cho2026}.

Oxley also compared the numbers of bases and circuits containing
a fixed independent set \cite[Theorem~1.3]{Oxley1983}.
Our second theorem compares bases avoiding a specified element with
all circuits of the original matroid.

\begin{theorem}\label{thm:deletion}
Let \(M\) be a nonempty simple binary matroid with no coloops, and \(e\in E(M)\). Then
\[
 b(M\backslash e)\ge d(M).
\]
Equality holds if and only if \(E(M)\) is a circuit, or \(M\) is a
parallel connection of a circuit and a triangle not containing \(e\).
In particular, the inequality is strict
when \(M\) is cosimple.
\end{theorem}

The deletion bound also holds when coloops are allowed. If \(e\) is
a coloop, then \(b(M\backslash e)=b(M)>d(M)\); otherwise, deleting
all coloops reduces the bound and its equality cases to
\autoref{thm:deletion}.
A one-element series extension preserves the number of circuits
and adds as many bases as avoid the extended element. Thus, we obtain
the following consequence.

\begin{corollary}\label{cor:growth}
Let \(N\) be a nonempty simple binary matroid with no coloops, and \(s\) be a nonnegative integer.
If \(M\) is obtained from \(N\) by \(s\) successive one-element
series extensions, then
\[
 b(M)\ge b(N)+s d(N).
\]
Equality holds if and only if \(s=0\), \(E(M)\) is a circuit, or
\(M\) is a parallel connection of a circuit and a triangle.
If \(N\) is cosimple and \(s\ge1\), the inequality is strict.
\end{corollary}

Oxley already used joint counts for two circuit sizes
\cite[proof of Lemma~2.11]{Oxley1983}.
We develop this approach by retaining the constraints imposed by the
simultaneous occurrence of circuits of different sizes as fundamental
circuits. Writing \(d_j(M)\) for the number of \(j\)-element circuits,
our central joint estimate is
\[
 \frac12d_{r(M)-1}(M)+d_{r(M)}(M)+\frac32d_{r(M)+1}(M)
 \le\frac{2b(M)}{r(M)+1}
\]
for simple cosimple binary matroids of rank at least five
(\autoref{prop:joint}). Its coefficients are adapted to the deletion
induction: a basis occurs in \(|E(M)|-r(M)\) single-element deletions,
whereas a \(j\)-element circuit occurs in \(|E(M)|-j\).
The joint estimate controls the resulting correction term
\(\sum_{j=3}^{r(M)+1}(j-r(M))d_j(M)\) and closes the induction.

To prove the joint estimate, we weight each occurrence of a fundamental
circuit \(C\) by \(1/b(M/C)\); over all bases, \(C\) then contributes
exactly \(|C|\). Rank-one and rank-two contractions control the weights
at a fixed basis. For each fixed \(r(M)\)-element circuit, the bases
at which it and a spanning circuit are both fundamental correspond
two-to-one to the edges of an induced subgraph of a folded cube;
see~\autoref{lem:cube-count}. Edge bounds of Ellis and Leader
\cite{EllisLeader2018}, combined with the local contraction bounds,
yield the required joint estimate.
For the deletion theorem, a separate basis-exchange count compensates
for the exceptional contributions of spanning circuits by additional
exchanges from distinct bases (\autoref{lem:exchange-count}).

The preliminary reductions are collected in \autoref{sec:prelim}.
\autoref{sec:weights} establishes the weighted and joint bounds.
\autoref{sec:proofs} combines deletion with reductions at series pairs
to prove \autoref{thm:rank}, then proves \autoref{thm:deletion}
and \autoref{cor:growth}.

\section{Preliminaries}\label{sec:prelim}

All matroids are finite. Terminology and notation follow
Oxley \cite{Oxley2011}.
We write \(E(M)\), \(\BB(M)\), \(\CC(M)\), and \(\cl_M\) for
the ground set, set of bases, set of circuits, and closure operator.
Deletion, contraction, restriction, and duality are denoted by
\(M\backslash X\), \(M/X\), \(M|X\), and \(M^*\), respectively.
We also write \(\CC^*(M)\coloneqq\CC(M^*)\) for the set of cocircuits.
For each integer \(j\ge1\), let \(d_j(M)\) be the number of \(j\)-element
circuits. For \(e\in E(M)\), let \(b_e(M)\) and \(d_e(M)\)
count the bases and circuits containing \(e\), respectively. We abbreviate
\(X-\{e\}\) to \(X-e\). A \emph{triangle} is a three-element circuit,
a \emph{triad} is a three-element cocircuit, and a \emph{series pair} is a
two-element cocircuit. A matroid is \emph{cosimple} if its dual is simple.
A \emph{parallel class} is a maximal set of nonloops any two
distinct members of which form a two-element circuit.

For \(B\in\BB(M)\) and \(e\in E(M)-B\), the
\emph{fundamental circuit} \(C(e,B)\) is the unique circuit
contained in \(B\cup\{e\}\). For \(f\in B\), the
\emph{fundamental cocircuit} \(C^*(f,B)\) is the unique cocircuit
contained in \((E(M)-B)\cup\{f\}\).
Basis exchange is characterized by
\cite[Section~1.2, Exercise~5 and Section~2.1, Exercise~10(b)]{Oxley2011}:
\begin{equation}\label{eq:exchange}
 (B-f)\cup\{e\}\in\BB(M)
 \quad\Longleftrightarrow\quad f\in C(e,B)
 \quad\Longleftrightarrow\quad e\in C^*(f,B).
\end{equation}
A matroid is \emph{binary} if it is representable over \(\FF\).
In a binary matroid \(M\), every circuit meets every cocircuit in an
even number of elements \cite[Theorem~9.1.2(ii)]{Oxley2011}:
\begin{equation}\label{eq:binary-orthogonality}
 |C\cap C^*|\equiv0\pmod2
 \qquad(C\in\CC(M),\ C^*\in\CC^*(M)).
\end{equation}
The symmetric difference of two circuits is a disjoint union
of circuits \cite[Theorem~9.1.2(iv)]{Oxley2011}:
\begin{equation}\label{eq:binary-symmetric-difference}
 C_1\triangle C_2=D_1\mathbin{\dot\cup}\cdots\mathbin{\dot\cup}D_k,
 \qquad C_1,C_2,D_1,\ldots,D_k\in\CC(M).
\end{equation}
The union may be empty. By duality,
\eqref{eq:binary-symmetric-difference} also holds for cocircuits.
Contraction preserves cosimplicity, since deletion preserves
simplicity in the dual. If \(M\) is cosimple and \(e\in E(M)\), then
\(e\) is not a coloop and \(M^*/e\) has no loops, since \(M^*\) is simple.
Thus, \(r(M\backslash e)=r(M)\), and
\(M\backslash e=(M^*/e)^*\) has no coloops.

When using a representation of a simple binary matroid \(M\),
we fix one of full row rank, with columns labelled by \(E(M)\).
In vector sums, each element denotes its column in \(\FF^{r(M)}\).
These columns are distinct and nonzero. For every circuit,
\begin{equation}\label{eq:circuit-column-sum}
 \sum_{e\in C}e=0\qquad(C\in\CC(M)).
\end{equation}
For a basis \(B\), put \(\sigma_B\coloneqq\sum_{x\in B}x\),
and use the columns of \(B\) as a coordinate basis.
Thus, \(\sigma_B\in E(M)\) means that this vector sum is itself
a column of the representation.
A circuit is \emph{spanning} if it has rank \(r(M)\), or
equivalently, size \(r(M)+1\).
If \(r(M)\ge2\), then
\begin{equation}\label{eq:spanning-bases}
 B\text{ lies in a spanning circuit}
 \quad\Longleftrightarrow\quad \sigma_B\in E(M),
\end{equation}
and that circuit is \(B\cup\{\sigma_B\}\).
Each spanning circuit contains \(r(M)+1\) bases,
and no basis belongs to two spanning circuits.
Consequently, for \(r(M)\ge2\),
\begin{equation}\label{eq:spanning-basis-count}
 \bigl|\{B\in\BB(M):\sigma_B\in E(M)\}\bigr|
 =(r(M)+1)d_{r(M)+1}(M).
\end{equation}
For every \(e\in E(F_7)\), we will also use the counts
\begin{equation}\label{eq:fano}
 b(F_7)=28,\quad d(F_7)=14,\quad
 b(F_7\backslash e)=16,\quad d(F_7\backslash e)=7.
\end{equation}

We refine the basis-extension argument in
\cite[proof of Theorem~1.3]{Oxley1983} as follows.
\begin{lemma}\label{lem:pointed}
If \(M\) is a simple binary matroid and \(e\in E(M)\), then
\[
 b_e(M)\ge d_e(M)+1.
\]
\end{lemma}
\begin{proof}
Fix a basis \(B_e\) containing \(e\). We construct an injection
from the circuits containing \(e\) to the bases containing \(e\)
other than \(B_e\).
The fundamental cocircuit satisfies
\(C^*(e,B_e)\cap B_e=\{e\}\).
For each circuit \(C\) containing \(e\), choose an element \(x\)
as follows. If \(C\cap B_e\ne\{e\}\), choose
\(x\in(C\cap B_e)-e\). If \(C\cap B_e=\{e\}\),
\eqref{eq:binary-orthogonality} makes
\((C\cap C^*(e,B_e))-e\) nonempty;
choose \(x\) in this set.
The set \(C-x\) is independent and spans \(C\).
Hence, \((C-x)\cup(B_e-C)\) spans \(M\). Extend \(C-x\)
within this set to a basis \(B_C\), fixing one such choice for each
\(C\). Then \(e\in B_C\) and \(x\notin B_C\).
If \(C\cap B_e\ne\{e\}\), then \(x\in B_e-B_C\), so \(B_C\ne B_e\).
If \(C\cap B_e=\{e\}\), then \(B_C\) contains the
\(|C|-2\ge1\) elements of \(C-\{e,x\}\) outside \(B_e\),
so again \(B_C\ne B_e\).

With these choices fixed for each \(C\), the map
\[
 \varphi:\{C\in\CC(M):e\in C\}\to
 \{B\in\BB(M):e\in B,\ B\ne B_e\},
\]
defined by \(\varphi(C)\coloneqq B_C\), is well defined by the construction.
We prove that \(\varphi\) is injective by distinguishing the
two cases in the construction by parity. Every element added when extending
\(C-x\) lies in \(B_e-C\), which is disjoint from \(C^*(e,B_e)\).
Thus, \(B_C\cap C^*(e,B_e)=(C-x)\cap C^*(e,B_e)\).
If \(C\cap B_e\ne\{e\}\), then \(x\notin C^*(e,B_e)\);
otherwise, \(x\in C^*(e,B_e)\).
By \eqref{eq:binary-orthogonality},
\begin{equation}\label{eq:pointed-parity}
 |B_C\cap C^*(e,B_e)|\equiv
 \begin{cases}
 0,&C\cap B_e\ne\{e\},\\
 1,&C\cap B_e=\{e\}
 \end{cases}
 \pmod2.
\end{equation}
Let \(C_1,C_2\) be circuits containing \(e\), and suppose
\(B_{C_1}=B_{C_2}\). By \eqref{eq:pointed-parity},
\(C_1\cap B_e=\{e\}\) if and only if \(C_2\cap B_e=\{e\}\).
If \(C_1\cap B_e\ne\{e\}\), then
\(C_1-B_e=B_{C_1}-B_e=B_{C_2}-B_e=C_2-B_e\), so
\(C_1\triangle C_2\subseteq B_e\).
If \(C_1\cap B_e=\{e\}\), then
\(C_i=(B_{C_i}-B_e)\cup\{e,x_i\}\) for \(i=1,2\), where \(x_i\)
is the element chosen for \(C_i\). Hence,
\(C_1\triangle C_2\subseteq\{x_1,x_2\}\).
Since \(B_e\) is a basis and \(M\) is simple,
\(C_1\triangle C_2\) is independent in both cases.
Thus, \eqref{eq:binary-symmetric-difference} implies
\(C_1\triangle C_2=\varnothing\), so \(\varphi\) is injective.
Consequently, \(d_e(M)\le b_e(M)-1\).
\end{proof}

To handle matroids with series pairs, we review the relevant
reductions and their basis and circuit counts.
A \emph{one-element series extension} \(M\) of \(N\) at a noncoloop
\(f\) replaces \(f\) by a series pair \(\{x,y\}\). Contracting \(x\)
and relabelling \(y\) as \(f\) recovers \(N\). The circuit set of \(M\) is
\begin{equation}\label{eq:series-circuits}
 \CC(M):=\CC(N\backslash f)\cup
 \{(C-f)\cup\{x,y\}:f\in C\in\CC(N)\}.
\end{equation}
For matroids \(N_1,N_2\) with ground sets meeting only at an
element \(f\) that is neither a loop nor a coloop of either,
their \emph{parallel connection} has circuit set
\begin{equation}\label{eq:parallel-circuits}
 \CC(P(N_1,N_2))\coloneqq\CC(N_1)\cup\CC(N_2)\cup
 \{(C_1\cup C_2)-f:f\in C_i\in\CC(N_i),\ i=1,2\}.
\end{equation}
These circuit descriptions appear in \cite[Section~7.1]{Oxley2011}.
In particular, a parallel connection of two circuits \(C_1,C_2\),
with \(C_1\cap C_2=\{f\}\), is binary
\cite[Proposition~7.1.24]{Oxley2011} and has exactly the three circuits
\(C_1\), \(C_2\), and \((C_1\cup C_2)-f\).
For any series pair \(\{x,y\}\), \eqref{eq:binary-orthogonality} gives
\(|\{x,y\}\cap C_i|\in\{0,2\}\) for \(i=1,2\).
Since \(C_1\cap C_2=\{f\}\), this forces
\(\{x,y\}\subseteq C_1-f\) or \(\{x,y\}\subseteq C_2-f\).

\begin{lemma}[{\cite[proof of Theorem~3.1, pp.~177--178]{Oxley1983}}]\label{lem:series}
Let \(M\) be a connected simple binary matroid of rank at least
three with a series pair. There are a connected simple binary minor
\(N\) with no coloops, where \(r(N)=r(M)-1\), and an element
\(f\in E(N)\), such that one of the following holds.
\begin{enumerate}
\item[\(\mathrm{(i)}\)] The matroid \(M\) is a one-element series
extension of \(N\) at \(f\), and
\[
 b(M)=b(N)+b(N\backslash f),\qquad d(M)=d(N).
\]
\item[\(\mathrm{(ii)}\)] The matroid \(M\) is the parallel connection
of \(N\) and a triangle \(\{f,x,y\}\), and
\[
 b(M)=2b(N)+b_f(N),\qquad d(M)=d(N)+d_f(N)+1.
\]
\end{enumerate}
\end{lemma}
The basis descriptions in \cite[Proposition~7.1.13(i),(ii)]{Oxley2011}
explain the basis counts.
In \autoref{lem:series}\(\mathrm{(i)}\), write \(\{x,y\}\) for the new series
pair. A basis \(B\) of \(N\) containing \(f\) corresponds to the basis
\((B-f)\cup\{x,y\}\) of \(M\); a basis avoiding \(f\) extends to the two
bases \(B\cup\{x\}\) and \(B\cup\{y\}\).
In \autoref{lem:series}\(\mathrm{(ii)}\), every basis \(B\) of \(N\) extends to
\(B\cup\{x\}\) and \(B\cup\{y\}\); if \(f\in B\), then
\((B-f)\cup\{x,y\}\) is also a basis. These descriptions list every basis
of \(M\) exactly once.
Since \(f\) is not a coloop of \(N\),
\(b(N)-b_f(N)=b(N\backslash f)\), which gives the stated basis counts.
The circuit identities follow from \eqref{eq:series-circuits}
and \eqref{eq:parallel-circuits}.
The series-extension identities in part \(\mathrm{(i)}\) also hold
without a connectedness assumption.
We also isolate the case in which deleting \(f\) removes every circuit.
If a matroid \(N\) has no coloops and \(f\in E(N)\) satisfies
\(d(N\backslash f)=0\), then \(E(N)-f\) is independent and
\(r(N\backslash f)=r(N)\), so it is a basis.
Hence, \(C(f,E(N)-f)\) is the only circuit of \(N\).
Any element outside this circuit would be a coloop. Thus,
\begin{equation}\label{eq:circuit-case}
 d(N\backslash f)=0\quad\Longrightarrow\quad E(N)\in\CC(N).
\end{equation}

\section{Weighted circuit bounds}\label{sec:weights}
Oxley \cite[Lemmas~2.2 and~2.7]{Oxley1983} proved the separate bounds
\[
(r(M)+1)d_{r(M)+1}(M)\le b(M)\qquad\text{and}\qquad r(M)d_{r(M)}(M)\le b(M)
\]
for every simple binary matroid \(M\).
Oxley's proof of Lemma~2.11 already counts circuits of different sizes
over the same bases \cite{Oxley1983}.
Here we develop joint estimates by assigning each fundamental circuit
\(C\) the weight \(1/b(M/C)\). This normalization makes each
\(j\)-element circuit contribute total weight \(j\) over all bases.
Rank-one and rank-two contractions then control the weights at each basis.
The further ingredient is an exact correspondence, for each fixed
\(r(M)\)-element circuit, between folded-cube edges and pairs of bases
at which that circuit and a spanning circuit are both fundamental
(\autoref{lem:cube-count}).
Combining these local and graph-theoretic estimates gives
\autoref{prop:joint}, the bound needed for the deletion induction.

\subsection{Counting and contraction bounds}

We first count pairs \((B,C)\) in which \(B\) is a basis of
\(M\) and \(C\) is a fundamental circuit with respect to \(B\).
Let \(X\subseteq E(M)\), and fix a basis \(B_X\) of \(M|X\).
The standard basis description of contraction
\cite[Corollary~3.1.8]{Oxley2011} is
\begin{equation}\label{eq:contraction-bases}
 \BB(M/X)=\{B\subseteq E(M)-X:B\cup B_X\in\BB(M)\}.
\end{equation}
For a circuit \(C\), each \(C-e\), with \(e\in C\), is a basis
of \(M|C\). By \eqref{eq:contraction-bases}, it extends to
exactly \(b(M/C)\) bases of \(M\). As \(e\) varies, these are
precisely the bases at which \(C\) is fundamental, each counted once.
Thus,
\begin{equation}\label{eq:fundamental-count}
 \bigl|\{B\in\BB(M):|B\cap C|=|C|-1\}\bigr|=|C|b(M/C).
\end{equation}
This is the contraction form of the double count in
\cite[proof of Theorem~1.1]{Oxley1983}.
Assign each pair \((B,C)\) weight \(1/b(M/C)\).
For a fixed basis \(B\), let \(\dloc{j}\) and \(\wloc{j}\)
be the number and total weight, respectively, of its
\(j\)-element fundamental circuits.
Explicitly,
\begin{align}
 \dloc{j}&\coloneqq\bigl|\{e\in E(M)-B:|C(e,B)|=j\}\bigr|,
 \notag\\
 \wloc{j}&\coloneqq\sum_{e\in E(M)-B,\,|C(e,B)|=j}
 \frac1{b(M/C(e,B))}.
 \label{eq:local-weight}
\end{align}
By \eqref{eq:fundamental-count}, a fixed \(j\)-element circuit
contributes total weight \(j\) over all bases.
Therefore, counting the same pairs by bases and by circuits,
with and without weights, we obtain
\begin{align}
 \sum_{B\in\BB(M)}\wloc{j}&=j d_j(M),
 \label{eq:weight-total}\\
 \sum_{B\in\BB(M)}\dloc{j}
 &=j\sum_{C\in\CC(M),\,|C|=j}b(M/C).
 \label{eq:count-total}
\end{align}
In particular, for a simple matroid,
\begin{equation*}
 d(M)=\sum_{B\in\BB(M)}\sum_{j=3}^{r(M)+1}\frac{\wloc{j}}j.
\end{equation*}

For a circuit \(C\), we have \(r(M/C)=r(M)-|C|+1\)
by the contraction rank formula \cite[Proposition~3.1.6]{Oxley2011}.
Thus, when \(|C|=r(M)\) or \(r(M)-1\), the contraction \(M/C\)
has rank one or two, respectively. If \(M\) is cosimple, so is
\(M/C\). In rank one, the set of nonloops is a cocircuit.
Thus, if \(M\) is cosimple and \(|C|=r(M)\), then \(b(M/C)\ge3\).
In rank two we use the following bounds, retaining any loops and
parallel elements of the contraction.

\begin{lemma}\label{lem:rank-two}
Let \(M\) be a cosimple binary matroid of rank two with \(t\)
nonloops.
\begin{enumerate}
\item[\(\mathrm{(i)}\)] We have
\[
 b(M)\ge\max\{8,3(t-3)\}.
\]
\item[\(\mathrm{(ii)}\)] If \(M\) has three nonempty parallel classes, then
\(b(M)\ge3t-7\).
\item[\(\mathrm{(iii)}\)] If \(M\) has three nonempty parallel classes,
each of size at least two, then \(b(M)\ge4t-12\).
\end{enumerate}
\end{lemma}
\begin{proof}
Four nonempty parallel classes would determine a \(U_{2,4}\)-restriction,
which is excluded by binarity \cite[Theorem~6.5.4]{Oxley2011}.
Thus, \(M\) has two or three nonempty parallel classes. Note that
a basis of \(M\) consists of two nonloops from distinct parallel
classes. If there are two parallel classes of sizes \(a,b\), then
\(t=a+b\). Each class is a cocircuit, so cosimplicity gives
\(a,b\ge3\). Hence,
\[
 b(M)=ab=3(t-3)+(a-3)(b-3)\ge\max\{9,3(t-3)\}.
\]
If there are three parallel classes, order their sizes as
\(1\le a\le b\le c\). Then \(t=a+b+c\) and
\(b(M)=ab+ac+bc\). The union of the classes of sizes \(a\) and \(b\)
is a cocircuit, so cosimplicity gives \(a+b\ge3\). Hence,
\(b,c\ge2\). Therefore,
\[
 b(M)-3t+7
 =(a-1)(b+c-3)+(b-2)(c-2)\ge0.
\]
Since \(t=a+b+c\ge5\), we also have
\(b(M)\ge3t-7\ge\max\{8,3(t-3)\}\).
If \(a,b,c\ge2\), then
\[
 b(M)-4t+12
 =(a-2)(b-2)+(a-2)(c-2)+(b-2)(c-2)\ge0.
\]
This completes the proof.
\end{proof}

\subsection{Weights at a fixed basis}

Throughout this subsection, \(M\) is a simple binary matroid,
with \(r(M)\ge4\), and \(B\in\BB(M)\).
We bound \(\wloc{r(M)}\) and \(\wloc{r(M)-1}\) by counting
nonloops and bases in the corresponding contractions.
We record the omitted basis elements so that these contraction
counts retain the dependence between the two circuit sizes.

An \(r(M)\)-element fundamental circuit omits exactly one
element \(x\) of \(B\). By \eqref{eq:circuit-column-sum}, its
element outside \(B\) is \(\sigma_B+x\). Thus, the set
\[
 X_B\coloneqq\{x\in B:\sigma_B+x\in E(M)\}
\]
records precisely
these circuits, with \(x\in X_B\) corresponding to
\[
 C_x\coloneqq(B-x)\cup\{\sigma_B+x\}.
\]
Similarly, an \((r(M)-1)\)-element fundamental circuit omits
a pair \(\{x,y\}\). Let \(G_B\) be the simple graph on \(B\)
whose edges are these omitted pairs: \(xy\in E(G_B)\) exactly
when \(\sigma_B+x+y\in E(M)\). The corresponding circuit is
\[
 C_{xy}\coloneqq(B-\{x,y\})\cup\{\sigma_B+x+y\}.
\]
Therefore, we have
\begin{equation}\label{eq:graph-counts}
 |X_B|=\dloc{r(M)},\qquad |E(G_B)|=\dloc{r(M)-1}.
\end{equation}

Write \(\delta_B(x)\) for the degree of \(x\) in \(G_B\). It counts the \((r(M)-1)\)-element fundamental circuits omitting \(x\).
For \(x\in X_B\), the circuit \(C_x\) is spanned by \(B-x\).
Thus, a representation of the rank-one contraction \(M/C_x\)
is obtained by retaining only the \(x\)-coordinate of each column outside \(C_x\).
The singleton sets \(\{x\}\), \(\{\sigma_B+z\}\) with
\(z\in X_B-x\), and \(\{\sigma_B+e+f\}\) with
\(ef\in E(G_B)\) and \(x\notin\{e,f\}\) are pairwise distinct.
The columns in these three families have, respectively, \(1\), \(r(M)-1\), and \(r(M)-2\)
nonzero \(B\)-coordinates, and the columns within each family are distinct.
Their elements all have \(x\)-coordinate \(1\), so these singleton sets
are bases of \(M/C_x\).
Counting these bases and using \(\delta_B(x)\le\dloc{r(M)-1}\), we obtain
\begin{equation}\label{eq:rank-one-selected}
 b(M/C_x)\ge
 \dloc{r(M)-1}+\dloc{r(M)}-\delta_B(x)
 \ge\dloc{r(M)}.
\end{equation}
If \(X_B\ne\varnothing\), each of the \(|X_B|\) summands in
\(\wloc{r(M)}\) is at most \(1/|X_B|\); if \(X_B=\varnothing\),
the sum is zero. Thus,
\begin{equation}\label{eq:general-rank-one-weight}
 \wloc{r(M)}\le1.
\end{equation}
If \(M\) is cosimple, then \(b(M/C_x)\ge3\) as well.

Next fix \(xy\in E(G_B)\). The set \(B-\{x,y\}\) is a basis
of \(M|C_{xy}\), so \(\{x,y\}\) is a basis of \(M/C_{xy}\)
by \eqref{eq:contraction-bases}.
Use a representation with column \((1,0)\) for \(x\)
and column \((0,1)\) for \(y\). Since \(C_{xy}\) is spanned
by \(B-\{x,y\}\), contraction retains just the \(x\)- and
\(y\)-coordinates of each remaining column.
Select \(x,y\), the columns \(\sigma_B+z\) with \(z\in X_B\),
and the columns \(\sigma_B+e+f\) with \(ef\in E(G_B)-xy\). These are
\(\dloc{r(M)-1}+\dloc{r(M)}+1\) nonloops. Each selected column is one of
\((0,1)\), \((1,0)\), and \((1,1)\).
Write \(\ind_{X_B}\) for the indicator of \(X_B\).
The respective numbers of selected elements are
\begin{equation}\label{eq:selected-classes}
\begin{split}
 &\delta_B(x)+\ind_{X_B}(x),\qquad
 \delta_B(y)+\ind_{X_B}(y),\\
 &\dloc{r(M)-1}+\dloc{r(M)}+1
 -\delta_B(x)-\delta_B(y)-\ind_{X_B}(x)-\ind_{X_B}(y).
\end{split}
\end{equation}
For example, the selected elements with column \((0,1)\) are \(y\),
the elements \(\sigma_B+x+z\) with \(xz\in E(G_B)\) and \(z\ne y\),
and \(\sigma_B+x\) when \(x\in X_B\).
The third number in \eqref{eq:selected-classes} equals
\(|X_B-\{x,y\}|\) plus the number of edges of \(G_B\) disjoint from \(xy\).
Since \(xy\in E(G_B)\), the first two numbers are positive.
The third is at least \(|X_B|-2\), so all three are positive if
\(\dloc{r(M)}\ge3\). If \(X_B=B\), the first two are at least two,
and the third is at least \(r(M)-2\ge2\).
Elements with the same nonzero column are parallel, and any two
selected elements from different parallel classes form a basis
of \(M/C_{xy}\). These counts concern only the selected
nonloops; any additional nonloops can only increase the number of bases.


The next lemma combines these counts with cosimplicity
to bound the weights at \(B\).
\begin{lemma}\label{lem:local}
For every basis \(B\) of a simple cosimple binary matroid \(M\)
with \(r(M)\ge4\), the following hold.
\begin{enumerate}
\item[\(\mathrm{(i)}\)] \(\wloc{r(M)}\le\min\{1,\dloc{r(M)}/3\}\).
\item[\(\mathrm{(ii)}\)]\(\wloc{r(M)-1}\le\max\{1/3,7/12-\dloc{r(M)}/8\}\).
\item[\(\mathrm{(iii)}\)] If \(\dloc{r(M)}>0\), then
\[
 \wloc{r(M)}\le
 \frac{\dloc{r(M)}^2+\sum_{x\in X_B}\delta_B(x)}
 {\dloc{r(M)}(\dloc{r(M)-1}+\dloc{r(M)})},
\]
where
\[
 \sum_{x\in X_B}\delta_B(x)
 \le\min\left\{2\dloc{r(M)-1},\,
 \dloc{r(M)-1}+\binom{\dloc{r(M)}}2\right\}.
\]
\item[\(\mathrm{(iv)}\)] If \(\dloc{r(M)}\ge2\), then
\[
 \wloc{r(M)}+3\wloc{r(M)-1}
 \le1+\frac2{\dloc{r(M)}}.
\]
\end{enumerate}
\end{lemma}
\begin{proof}
Each \(x\in X_B\) contributes \(1/b(M/C_x)\) to \(\wloc{r(M)}\).
By \eqref{eq:rank-one-selected} and cosimplicity,
\(b(M/C_x)\ge\max\{3,\dloc{r(M)}\}\). Hence,
\[
 \wloc{r(M)}\le
 \frac{\dloc{r(M)}}{\max\{3,\dloc{r(M)}\}}
 =\min\{1,\dloc{r(M)}/3\},
\]
proving \(\mathrm{(i)}\).

For each \(xy\in E(G_B)\), the contraction \(M/C_{xy}\)
is cosimple and binary of rank two. By \eqref{eq:selected-classes},
it has at least \(\dloc{r(M)-1}+\dloc{r(M)}+1\) nonloops.
Applying \autoref{lem:rank-two}\(\mathrm{(i)}\) and
expanding \eqref{eq:local-weight}, we obtain
\begin{equation}\label{eq:rank-two-selected}
 \wloc{r(M)-1}=\sum_{xy\in E(G_B)}\frac1{b(M/C_{xy})}\le\frac{\dloc{r(M)-1}}{\max\{8,3(\dloc{r(M)-1}+\dloc{r(M)}-2)\}}.
\end{equation}
Decompose the numerator in \eqref{eq:rank-two-selected} as
\((\dloc{r(M)-1}+\dloc{r(M)}-2)+(2-\dloc{r(M)})\).
The denominator is positive, so a negative numerator term gives
a nonpositive quotient. Therefore, the two lower bounds on the denominator
bound the quotients by \(1/3\) and
\(\max\{0,2-\dloc{r(M)}\}/8\), respectively. Adding the bounds yields
\[
 \wloc{r(M)-1}\le\frac13+\frac{\max\{0,2-\dloc{r(M)}\}}8
 =\max\Bigl\{\frac13,\frac7{12}-\frac{\dloc{r(M)}}8\Bigr\},
\]
which proves \(\mathrm{(ii)}\).

For \(\mathrm{(iii)}\), suppose \(\dloc{r(M)}>0\).
Expanding \eqref{eq:local-weight} and using \eqref{eq:rank-one-selected}, we obtain
\begin{align*}
 (\dloc{r(M)-1}+\dloc{r(M)})\wloc{r(M)}&=\sum_{x\in X_B}\frac{\dloc{r(M)-1}+\dloc{r(M)}}{b(M/C_x)}\\
 &\le\sum_{x\in X_B}\Bigl(1+\frac{\delta_B(x)}{b(M/C_x)}\Bigr)\\
 &\le\dloc{r(M)}+\frac{\sum_{x\in X_B}\delta_B(x)}{\dloc{r(M)}}.
\end{align*}
The last step uses \(b(M/C_x)\ge\dloc{r(M)}\) from
\eqref{eq:rank-one-selected}.
Here \(|X_B|=\dloc{r(M)}\) by \eqref{eq:graph-counts}.
Dividing by \(\dloc{r(M)-1}+\dloc{r(M)}\) proves the first inequality in
\(\mathrm{(iii)}\).
Each edge of \(G_B\) contributes at most two to
\(\sum_{x\in X_B}\delta_B(x)\), proving the first degree-sum bound in
\(\mathrm{(iii)}\).
At most \(\binom{|X_B|}2=\binom{\dloc{r(M)}}2\) edges have both ends in
\(X_B\) and contribute twice; every other edge contributes at most once.
This proves the second degree-sum bound in \(\mathrm{(iii)}\).

For \(\mathrm{(iv)}\), the case \(\dloc{r(M)-1}=0\)
follows from \(\mathrm{(i)}\). Otherwise, the denominators below
are positive, since \(\dloc{r(M)-1}>0\) and \(\dloc{r(M)}\ge2\).
Combining \(\mathrm{(iii)}\) with \eqref{eq:rank-two-selected}, we derive
\[
 \wloc{r(M)}\le\frac{\dloc{r(M)}^2+2\dloc{r(M)-1}}
 {\dloc{r(M)}(\dloc{r(M)-1}+\dloc{r(M)})},
\]
\[
 3\wloc{r(M)-1}\le\frac{\dloc{r(M)-1}}{\dloc{r(M)-1}+\dloc{r(M)}-2}\le\frac{\dloc{r(M)-1}+2}{\dloc{r(M)-1}+\dloc{r(M)}}.
\]
Clearing the positive denominators reduces the last inequality to
\(2(\dloc{r(M)}-2)\ge0\).
The two upper bounds sum to \(1+2/\dloc{r(M)}\), proving \(\mathrm{(iv)}\). This completes the proof.
\end{proof}

\subsection{A folded-cube estimate}

We next relate circuits of size \(r(M)\) to spanning circuits
by counting the bases at which both are fundamental.
For a fixed \(r(M)\)-element circuit, this count will be twice
the number of edges induced by a suitable vertex subset of a
folded \(r(M)\)-cube. We first prove the required edge bound.

For an integer \(\ell\ge4\), the \emph{folded \(\ell\)-cube}, denoted by \(\square_\ell\), is
obtained from the ordinary \((\ell-1)\)-cube by adding an edge
between each pair of antipodal vertices. Its vertices are the
vectors in \(\FF^{\ell-1}\): edges of the ordinary \((\ell-1)\)-cube join
vectors differing in exactly one coordinate, and antipodal edges join
vectors differing in every coordinate.
For \(U\subseteq V(\square_\ell)\), let
\(\varepsilon_{\square_\ell}(U)\) be the number of edges with
both endpoints in \(U\), that is, the number of edges of
the induced subgraph \(\square_\ell[U]\).

Following Ellis and Leader \cite{EllisLeader2018}, for each integer
\(t\ge0\), let \(F(t)\) count the edges induced by the first \(t\)
vertices of an ordinary cube with at least \(t\) vertices, ordered
lexicographically by their coordinates, with \(0<1\).
Coordinate reversal identifies this order with that used by Ellis and Leader
and preserves the induced edge count.
Increasing the ambient dimension only adds initial zero coordinates to these vertices,
so their induced edge count is unchanged; also \(F(0)=0\).
Since \(U\subseteq V(\square_\ell)=\FF^{\ell-1}\), we have
\(|U|\le2^{\ell-1}\), so \(F(|U|)\) can be computed in the
ordinary \((\ell-1)\)-cube.
In dimension \(\ell-1\), Theorem~3 of \cite{EllisLeader2018} applies for
\(0\le|U|\le2^{\ell-2}\), and equation~(7) applies for
\(2^{\ell-2}\le|U|\le2^{\ell-1}\). Neither requires \(U\) to be antipodal.
The quantity bounded in both results is \(2\varepsilon_{\square_\ell}(U)\):
every edge of the ordinary \((\ell-1)\)-cube induced by \(U\),
and every antipodal pair in \(U\), is counted twice.
Dividing their bounds by two, we obtain
\begin{equation}\label{eq:cube-initial}
 \varepsilon_{\square_\ell}(U)
 \le F(|U|)+\max\{0,|U|-2^{\ell-2}\}.
\end{equation}
We derive the following quadratic bounds from \eqref{eq:cube-initial}.

\begin{lemma}\label{lem:cube}
For \(\ell\ge4\) and \(U\subseteq V(\square_\ell)\), we have
\[
 \varepsilon_{\square_\ell}(U)\le
 \begin{cases}
 |U|^2/4,&\ell=4,\\
 |U|/2+|U|^2/8,&\ell=5,\\
 |U|+|U|^2/16,&\ell\ge6.
 \end{cases}
\]
\end{lemma}
\begin{proof}
We first bound the ordinary-cube term \(F(t)\), then absorb
the additional term in \eqref{eq:cube-initial}.
For \(k\in\{1,2,3\}\), set
\(G_k(t)\coloneqq(k-1)t/2+t^2/2^{k+1}\). We prove by induction on \(k\) that
\begin{equation}\label{eq:ordinary-cube}
 F(t)\le G_k(t)\qquad(0\le t\le2^{\ell-1}).
\end{equation}
For \(k=1\), the subgraph defining \(F(t)\) is bipartite.
Hence, \(F(t)\le t^2/4=G_1(t)\).

Let \(k\in\{2,3\}\), and assume that \eqref{eq:ordinary-cube} holds
with \(k-1\) in place of \(k\) for every integer \(0\le t\le2^{\ell-1}\).
Write \(t=2^k a+s\), where \(a,s\) are integers with \(a\ge0\) and \(0\le s<2^k\).
Since \(\ell-1\ge k\), partition the first \(t\) vertices of the ordinary
\((\ell-1)\)-cube according to their first \(\ell-1-k\) coordinates.
This yields \(a\) vertex sets, each inducing an ordinary \(k\)-cube,
and a possibly empty remaining set of \(s\) vertices.
All vertices in the remaining set have the same first \(\ell-1-k\) coordinates.
Deleting these common coordinates preserves adjacency and leaves the first
\(s\) vertices of the ordinary \(k\)-cube in lexicographic order, with \(0<1\).
Thus, by the definition of \(F\), the remaining set induces exactly \(F(s)\) edges.
Each ordinary \(k\)-cube has \(k2^{k-1}\) edges.
An edge joining distinct sets leaves the last \(k\) coordinates unchanged,
so each vertex has at most one neighbour in any other set.
Thus, at most \(2^k\) edges join any two of these ordinary \(k\)-cubes,
and at most \(s\) edges join each ordinary \(k\)-cube to the remaining set.
Since \(s<2^k\le2^{\ell-1}\), the induction hypothesis yields
\(F(s)\le G_{k-1}(s)\).
Counting the induced edges, we obtain
\[
 F(t)\le k2^{k-1}a+2^k\binom a2+as+G_{k-1}(s)=G_k(t)-\frac{s(2^k-s)}{2^{k+1}}\le G_k(t).
\]
This proves \eqref{eq:ordinary-cube}.

Choose \(k=1\) if \(\ell=4\), \(k=2\) if \(\ell=5\), and
\(k=3\) otherwise. Then \(\ell\ge k+3\).
For \(|U|\le2^{\ell-2}\), \eqref{eq:cube-initial} and
\eqref{eq:ordinary-cube} prove the required bound.
Otherwise, write \(|U|=2^{\ell-2}+s\), where \(1\le s\le2^{\ell-2}\).
In the ordinary \((\ell-1)\)-cube, the first \(2^{\ell-2}\) vertices
are those with first coordinate \(0\) and induce an ordinary \((\ell-2)\)-cube.
The next \(s\) vertices have first coordinate \(1\), induce \(F(s)\) edges,
and each has exactly one neighbour in the first subcube. Hence,
\(F(2^{\ell-2}+s)=F(2^{\ell-2})+s+F(s)\).
Using \eqref{eq:cube-initial} and \eqref{eq:ordinary-cube}, we obtain
\[
\varepsilon_{\square_\ell}(U)\le G_k(2^{\ell-2})+2s+G_k(s)
=G_k(|U|)+(2-2^{\ell-k-2})s\le G_k(|U|).
\]
The last inequality uses \(\ell\ge k+3\) and \(s\ge0\).
This completes the proof.
\end{proof}


We now apply this edge bound to count the bases at which an
\(r(M)\)-element circuit and a spanning circuit are both fundamental.


\begin{lemma}\label{lem:cube-count}
Let \(M\) be a simple binary matroid with \(r(M)\ge4\).
Then the following hold.
\begin{enumerate}
\item[\(\mathrm{(i)}\)] If \(r(M)=4\), then
\[
 4d_4(M)+5d_5(M)
 \le b(M)+\frac18\sum_{B\in\BB(M)}\dloc4.
\]
\item[\(\mathrm{(ii)}\)] If \(r(M)=5\), then
\[
 4d_5(M)+6d_6(M)
 \le b(M)+\frac1{20}\sum_{B\in\BB(M)}\dloc5.
\]
\item[\(\mathrm{(iii)}\)] If \(r(M)\ge6\), then
\[
 (r(M)-2)d_{r(M)}(M)+(r(M)+1)d_{r(M)+1}(M)
 \le b(M)+\frac1{8r(M)}\sum_{B\in\BB(M)}\dloc{r(M)}.
\]
\end{enumerate}
\end{lemma}

\begin{proof}
We separate the bases according to whether they lie in a spanning circuit.
By \eqref{eq:spanning-basis-count} and \eqref{eq:general-rank-one-weight},
\begin{equation}\label{eq:basis}
 \sum_{B\in\BB(M),\,\sigma_B\notin E(M)}\wloc{r(M)}
 \le b(M)-(r(M)+1)d_{r(M)+1}(M).
\end{equation}
We express the complementary sum as a folded \(r(M)\)-cube edge count.

Fix an \(r(M)\)-element circuit \(C\), and let
\(U_C=E(M)-\cl_M(C)\).
The nonloops of the rank-one contraction \(M/C\) are precisely
the elements of \(U_C\). Hence,
\begin{equation}\label{eq:outside-hyperplane}
 |U_C|=b(M/C).
\end{equation}
Since \(r(M|C)=r(M)-1\), the space \(\spanof(C)\) is a hyperplane in \(\FF^{r(M)}\).
On its other coset \(\FF^{r(M)}-\spanof(C)\), join two vectors
when their sum belongs to \(C\). Any \(r(M)-1\) columns of \(C\) determine the
directions of an ordinary \((r(M)-1)\)-cube; by \eqref{eq:circuit-column-sum},
the remaining column is their sum and specifies the antipodal edges.
Thus, this graph is a copy of \(\square_{r(M)}\). Under this identification,
\(U_C\subseteq V(\square_{r(M)})\) and
\(0<|U_C|\le2^{r(M)-1}\).

For an edge \(\{e,f\}\) of \(\square_{r(M)}[U_C]\), put
\(z=e+f\in C\). Since \(C-z\) is a basis of \(\spanof(C)\), this edge gives the two bases
\[
 (C-z)\cup\{e\},\qquad (C-z)\cup\{f\},
\]
at which \(C\) is fundamental. By \eqref{eq:circuit-column-sum},
their column sums are \(f\) and \(e\), respectively, so adjoining
the other endpoint produces a spanning circuit.
Conversely, if \(C\) is fundamental at
\(B\) and \(\sigma_B\in E(M)\), let
\(C-B=\{z\}\) and \(B-C=\{e\}\).
Since \(C-z\) spans \(C\), the basis property gives \(e\in U_C\).
By \eqref{eq:circuit-column-sum}, \(\sigma_B=e+z\in U_C\) as well.
Thus, \(B\) uniquely determines the edge \(\{e,\sigma_B\}\) of
\(\square_{r(M)}[U_C]\) and its endpoint \(e\).
Therefore, each edge of \(\square_{r(M)}[U_C]\) corresponds to exactly
two bases \(B\) at which \(C\) is fundamental and \(\sigma_B\in E(M)\).
By \eqref{eq:local-weight}, we obtain
\begin{align}
 \sum_{B\in\BB(M),\,\sigma_B\in E(M)}\wloc{r(M)}&=\sum_{C\in\CC(M),\,|C|=r(M)}
 \frac{\bigl|\{B\in\BB(M):|C-B|=1,\ \sigma_B\in E(M)\}\bigr|}{b(M/C)}\notag\\
 &=2\sum_{C\in\CC(M),\,|C|=r(M)}
 \frac{\varepsilon_{\square_{r(M)}}(U_C)}{b(M/C)}.
 \label{eq:cube-weight-count}
\end{align}
By \eqref{eq:spanning-bases}, the bases lying in no spanning circuit
are exactly those with \(\sigma_B\notin E(M)\).
Splitting the sum in \eqref{eq:weight-total} and applying
\eqref{eq:basis} and \eqref{eq:cube-weight-count}, we obtain
\begin{align*}
 r(M)d_{r(M)}(M)
 &=\sum_{\substack{B\in\BB(M)\\\sigma_B\notin E(M)}}\wloc{r(M)}
 +\sum_{\substack{B\in\BB(M)\\\sigma_B\in E(M)}}\wloc{r(M)}\\
 &\le b(M)-(r(M)+1)d_{r(M)+1}(M)+2\sum_{C\in\CC(M),\,|C|=r(M)}
 \frac{\varepsilon_{\square_{r(M)}}(U_C)}{b(M/C)}.
\end{align*}
Rearranging yields
\begin{equation}\label{eq:cube-master}
 r(M)d_{r(M)}(M)+(r(M)+1)d_{r(M)+1}(M)
 \le b(M)+2\sum_{C\in\CC(M),\,|C|=r(M)}
 \frac{\varepsilon_{\square_{r(M)}}(U_C)}{b(M/C)}.
\end{equation}
With \(\ell=r(M)\ge4\) and \(U=U_C\subseteq V(\square_{r(M)})\),
\autoref{lem:cube} and \eqref{eq:outside-hyperplane} imply
\[
 \frac{2\varepsilon_{\square_{r(M)}}(U_C)}{b(M/C)}
 \le\begin{cases}
 b(M/C)/2,&r(M)=4,\\
 1+b(M/C)/4,&r(M)=5,\\
 2+b(M/C)/8,&r(M)\ge6.
 \end{cases}
\]
By \eqref{eq:count-total}, we have
\[
 \sum_{C\in\CC(M),\,|C|=r(M)}b(M/C)
 =\frac1{r(M)}\sum_{B\in\BB(M)}\dloc{r(M)}.
\]
Substituting this identity and the preceding bounds in
\eqref{eq:cube-master} proves the three assertions.
\end{proof}


\subsection{The joint bounds}

We now combine the local counts retained in \autoref{lem:cube-count}
with the circuit-weight bounds in \autoref{lem:local}, at each basis.
In rank four, every circuit has size three, four, or five,
so these estimates bound \(d(M)\) directly.

\begin{proposition}\label{prop:rank-four}
Let \(M\) be a simple cosimple binary matroid of rank four. Then
\[
 d(M)<\frac25 b(M).
\]
\end{proposition}
\begin{proof}
To bound \(d(M)=d_3(M)+d_4(M)+d_5(M)\), we eliminate \(d_5(M)\)
using \autoref{lem:cube-count}\(\mathrm{(i)}\) and apply \eqref{eq:weight-total}:
\begin{align}
 d(M)&\le\frac15b(M)+d_3(M)+\frac15d_4(M)
       +\frac1{40}\sum_{B\in\BB(M)}\dloc4\notag\\
 &=\frac15b(M)+\sum_{B\in\BB(M)}
 \Bigl(\frac{\wloc3}3+\frac{\wloc4}{20}+\frac{\dloc4}{40}\Bigr).
 \label{eq:rank-four-target}
\end{align}
By \eqref{eq:rank-four-target}, it suffices to show, for each \(B\), that
\begin{equation}\label{eq:rank-four-local-target}
 20\wloc3+3\wloc4+\frac32\dloc4<12.
\end{equation}
By \eqref{eq:graph-counts},
\[
 0\le\dloc3\le6,\qquad 0\le\dloc4\le4.
\]
If \(\dloc4\le2\), \autoref{lem:local}\(\mathrm{(i),(ii)}\) imply
\(\wloc4\le\dloc4/3\) and
\(\wloc3\le7/12-\dloc4/8\). Hence,
\[
 20\wloc3+3\wloc4+\frac32\dloc4
 \le20\Bigl(\wloc3+\frac{\dloc4}8\Bigr)\le\frac{35}3<12.
\]
If \(\dloc4=3\), all three multiplicities in
\eqref{eq:selected-classes} are positive.
Each \(M/C_{xy}\), with \(xy\in E(G_B)\), has at least
\(\dloc3+4\) nonloops.
Using the degree-sum bound \(\sum_{x\in X_B}\delta_B(x)\le\dloc3+3\)
in \autoref{lem:local}\(\mathrm{(iii)}\), together with
\autoref{lem:rank-two}\(\mathrm{(ii)}\), we obtain
\[
 \wloc3\le\frac{\dloc3}{3\dloc3+5},\qquad
 \wloc4\le\frac13+\frac3{\dloc3+3}.
\]
Consequently, the left-hand side of \eqref{eq:rank-four-local-target}
is at most
\[
 \frac{11}2+\frac{20\dloc3}{3\dloc3+5}+\frac9{\dloc3+3}
 =12-\frac{(15-\dloc3)(\dloc3+7)}
 {2(3\dloc3+5)(\dloc3+3)}<12.
\]
If \(\dloc4=4\), all three multiplicities in
\eqref{eq:selected-classes} are at least two.
Each \(M/C_{xy}\), with \(xy\in E(G_B)\), has at least
\(\dloc3+5\) nonloops. Moreover,
\(X_B=B\), so \(\sum_{x\in X_B}\delta_B(x)=2\dloc3\).
By \autoref{lem:rank-two}\(\mathrm{(iii)}\) and
\autoref{lem:local}\(\mathrm{(iii)}\),
\[
 \wloc3\le\frac{\dloc3}{4\dloc3+8},\qquad
 \wloc4\le1-\frac{\dloc3}{2(\dloc3+4)}.
\]
Therefore, the left-hand side of \eqref{eq:rank-four-local-target} is at most
\[
 9+\frac{5\dloc3}{\dloc3+2}-\frac{3\dloc3}{2(\dloc3+4)}
 =12-\frac{(8-\dloc3)(\dloc3+6)}
 {2(\dloc3+2)(\dloc3+4)}<12.
\]
Now \eqref{eq:rank-four-target} proves the assertion.
\end{proof}

For higher ranks, the deletion argument in \autoref{sec:proofs}
requires a joint bound for the three largest possible circuit sizes.
The coefficients below are chosen to control the correction term
\(\sum_{j=3}^{r(M)+1}(j-r(M))d_j(M)\) arising when the induction
bounds are summed over all single-element deletions.
We establish that bound next.
\begin{proposition}\label{prop:joint}
Let \(M\) be a simple cosimple binary matroid with \(r(M)\ge5\). Then
\[
 \frac12d_{r(M)-1}(M)+d_{r(M)}(M)+\frac32d_{r(M)+1}(M)
 \le\frac{2b(M)}{r(M)+1}.
\]
\end{proposition}
\begin{proof}
We scale the bound in \autoref{lem:cube-count} to match the coefficient of
\(d_{r(M)+1}(M)\), then bound the remaining terms at each basis.

\emph{Rank five.}
Multiply the inequality in \autoref{lem:cube-count}\(\mathrm{(ii)}\)
by \(3/2\), then add \(3d_4(M)\). Using \eqref{eq:weight-total}, we obtain
\begin{equation}
 6\Bigl(\frac12d_4(M)+d_5(M)+\frac32d_6(M)\Bigr)\le\frac32b(M)+\sum_{B\in\BB(M)}
 \Bigl(\frac34\wloc4+\frac3{40}\dloc5\Bigr).
 \label{eq:rank-five-target}
\end{equation}
By \eqref{eq:graph-counts}, \(0\le\dloc4\le10\) and \(0\le\dloc5\le5\).
We prove that each summand is at most \(1/2\).
If \(\dloc4=0\), it is at most \(3/8\).
Suppose \(\dloc4>0\).
For \(\dloc5\le3\), we estimate the summand by \autoref{lem:local}\(\mathrm{(ii)}\):
\[
 \frac34\wloc4+\frac3{40}\dloc5
 \le\max\Bigl\{\frac14+\frac3{40}\dloc5,
 \frac7{16}-\frac3{160}\dloc5\Bigr\}
 \le\frac{19}{40}<\frac12.
\]
For \(\dloc5=4\), every \(M/C_{xy}\), with \(xy\in E(G_B)\),
has three nonempty parallel classes and at least \(\dloc4+5\) nonloops
by \eqref{eq:selected-classes}.
Since \(\dloc4\le10\),
\autoref{lem:rank-two}\(\mathrm{(ii)}\) shows that
every four-element fundamental circuit \(C\) satisfies
\[
 b(M/C)\ge3\dloc4+8\ge\frac{15}4\dloc4.
\]
Thus, \(\wloc4\le4/15\), so
\(3\wloc4/4+3\dloc5/40\le1/2\). It remains to prove \(\wloc4\le1/6\) when \(\dloc5=5\).
Fix \(xy\in E(G_B)\), and let \(a,b,c\)
count the selected elements in the three parallel classes of
\(M/C_{xy}\), in the order displayed in \eqref{eq:selected-classes}.
Here \(X_B=B\), so \(a+b+c=\dloc4+6\),
\(a,b\le5\), and \(c\le3+\binom32=6\).
Counting bases formed by two selected elements in different
parallel classes yields
\[
 b(M/C_{xy})\ge ab+ac+bc
 =6\dloc4+\frac{abc+(6-a)(6-b)(6-c)}6
 \ge6\dloc4.
\]
It follows that \(\wloc4\le1/6\), and
\(3\wloc4/4+3\dloc5/40\le1/2\).
Now \eqref{eq:rank-five-target} proves the assertion in rank five.

\emph{Rank at least six.}
Scale the bound in \autoref{lem:cube-count}\(\mathrm{(iii)}\)
by \(3/2\), and add
\[
 \frac{r(M)+1}2d_{r(M)-1}(M)+\frac{8-r(M)}2d_{r(M)}(M)
\]
to both sides.
Using \eqref{eq:weight-total}, we rewrite this inequality as
\begin{align}
 &(r(M)+1)\Bigl(\frac12d_{r(M)-1}(M)+d_{r(M)}(M)+\frac32d_{r(M)+1}(M)\Bigr)\le\frac32b(M)+\notag\\
 &\sum_{B\in\BB(M)}\biggl(\frac{r(M)+1}{2(r(M)-1)}\wloc{r(M)-1}+\frac{8-r(M)}{2r(M)}\wloc{r(M)}+\frac3{16r(M)}\dloc{r(M)}\biggr).
 \label{eq:high-rank-target}
\end{align}
Since \((r(M)+1)/(2(r(M)-1))\le3/4\) and
\((8-r(M))/(2r(M))\le1/6\), the nonnegativity of the weights
shows that it suffices to prove
\begin{equation}\label{eq:high-rank-local-target}
 \frac34\wloc{r(M)-1}+\frac16\wloc{r(M)}
 +\frac3{16r(M)}\dloc{r(M)}\le\frac12.
\end{equation}
If \(\dloc{r(M)}\le1\), \autoref{lem:local}\(\mathrm{(i),(ii)}\)
and \(r(M)\ge6\)
bound the left-hand side of
\eqref{eq:high-rank-local-target} by
\[
 \frac34\Bigl(\frac7{12}-\frac{\dloc{r(M)}}8\Bigr)
 +\frac{\dloc{r(M)}}{18}+\frac{\dloc{r(M)}}{32}
 =\frac7{16}-\frac{\dloc{r(M)}}{144}<\frac12.
\]
If \(2\le\dloc{r(M)}\le r(M)\), then
\autoref{lem:local}\(\mathrm{(ii)}\) gives
\(\wloc{r(M)-1}\le1/3\).
Multiply this bound by \(1/4\) and
\autoref{lem:local}\(\mathrm{(iv)}\) by \(1/6\), and add.
Thus, the left-hand side of \eqref{eq:high-rank-local-target} is at most
\[
 \frac14+\frac1{3\dloc{r(M)}}+\frac{3\dloc{r(M)}}{16r(M)}
 \le\frac7{16}+\frac1{3r(M)}<\frac12.
\]
The first inequality follows from
\begin{align*}
 &\frac7{16}+\frac1{3r(M)}-\Bigl(\frac14+\frac1{3\dloc{r(M)}}+\frac{3\dloc{r(M)}}{16r(M)}\Bigr)\\
 &=\frac{r(M)-\dloc{r(M)}}{r(M)}\Bigl(\frac3{16}-\frac1{3\dloc{r(M)}}\Bigr)\ge0.
\end{align*}
Equations~\eqref{eq:high-rank-target} and
\eqref{eq:high-rank-local-target} complete the proof.
\end{proof}

\section{Proofs of the main results}\label{sec:proofs}

We first prove \autoref{thm:rank} by induction, using the bounds from
\autoref{sec:weights} for the cosimple case and \autoref{lem:series} for
the series-pair case. For \autoref{thm:deletion}, the rank bound and a
basis-exchange estimate handle the cosimple case; induction and the same
series-pair reductions handle the remaining cases. Finally, we iterate
the series-extension formula to prove \autoref{cor:growth}.

\begin{proof}[Proof of \autoref{thm:rank}]
We use strong induction on \(|E(M)|\), including the equality classification.
The hypotheses imply \(r(M)\ge2\) and \(d(M)\ge1\).
For \(r(M)\le3\), \eqref{eq:oxley} implies
\(2b(M)\ge4d(M)\ge(r(M)+1)d(M)\).
The second inequality is strict when \(r(M)=2\); when \(r(M)=3\),
equality in the chain holds precisely for \(M\cong F_7\).
Assume henceforth that \(r(M)\ge4\).

\emph{Disconnected matroids.}
Write \(M=M_1\oplus M_2\), with both summands nonempty.
The circuit counts and ranks add, while the basis counts multiply
\cite[(4.2.12)--(4.2.14)]{Oxley2011}.
Each summand has no coloops, so \(r(M_i)\ge2\).
By \eqref{eq:exchange} and the absence of coloops, every element of a
fixed basis can be exchanged for an element outside it, producing
\(r(M_i)\) further distinct bases.
Thus, \(b(M_i)\ge r(M_i)+1\).
If both ranks are at least three, the induction hypothesis and
\(b(M_i)\ge r(M_i)+1\) imply
\[
 \frac{d(M)}{b(M)}
 =\frac{d(M_1)+d(M_2)}{b(M_1)b(M_2)}
 \le\frac4{(r(M_1)+1)(r(M_2)+1)}
 <\frac2{r(M)+1}.
\]
If one summand has rank two, relabel the summands so that
\(r(M_1)=2\). Then \(M_1\cong U_{2,3}\).
Adding \(2b(M_2)\ge(r(M_2)+1)d(M_2)\) to
\(b(M_2)\ge r(M_2)+1\) yields
\[
 \frac{d(M)}{b(M)}
 =\frac{1+d(M_2)}{3b(M_2)}
 \le\frac1{r(M_2)+1}
 <\frac2{r(M_2)+3}=\frac2{r(M)+1}.
\]

\emph{Connected cosimple matroids.}
Rank four is covered by \autoref{prop:rank-four}.
Suppose \(r(M)\ge5\).
Since \(M\) is cosimple, every deletion has rank \(r(M)\)
and no coloops.
No deletion is isomorphic to \(F_7\), since
\(r(M\backslash e)=r(M)\ge5\). Thus, the induction hypothesis gives
\[
 2b(M\backslash e)>(r(M\backslash e)+1)d(M\backslash e)
 =(r(M)+1)d(M\backslash e)\qquad(e\in E(M)).
\]
As in \cite[proof of Theorem~3.1, pp.~176--177]{Oxley1983},
we sum the resulting positive differences over all single-element deletions.
The resulting circuit-size correction will be controlled by
\autoref{prop:joint}.
Each basis occurs in \(|E(M)|-r(M)\) deletions, and each
\(j\)-element circuit occurs in \(|E(M)|-j\) deletions.
Therefore,
\begin{align*}
 0&<\sum_{e\in E(M)}
 \bigl(2b(M\backslash e)-(r(M)+1)d(M\backslash e)\bigr)\\
 &=2(|E(M)|-r(M))b(M)
 -(r(M)+1)\sum_{j=3}^{r(M)+1}(|E(M)|-j)d_j(M)\\
 &=(|E(M)|-r(M))\bigl(2b(M)-(r(M)+1)d(M)\bigr)+(r(M)+1)\sum_{j=3}^{r(M)+1}(j-r(M))d_j(M).
\end{align*}
To control the last sum, use \(j-r(M)\le-2\) for
\(3\le j\le r(M)-2\) and apply \autoref{prop:joint}:
\begin{align*}
 \sum_{j=3}^{r(M)+1}(j-r(M))d_j(M)
 &\le -2\sum_{j=3}^{r(M)-2}d_j(M)-d_{r(M)-1}(M)+d_{r(M)+1}(M)\\
 &=2\Bigl(\frac12d_{r(M)-1}(M)+d_{r(M)}(M)
          +\frac32d_{r(M)+1}(M)-d(M)\Bigr)\\
 &\le\frac{4b(M)}{r(M)+1}-2d(M).
\end{align*}
Substituting the inequality into the expansion of the sum over deletions, we obtain
\[
 0<\sum_{e\in E(M)}\bigl(2b(M\backslash e)-(r(M)+1)d(M\backslash e)\bigr)\le(|E(M)|-r(M)+2)\bigl(2b(M)-(r(M)+1)d(M)\bigr).
\]
Since \(|E(M)|-r(M)+2>0\), the desired strict bound follows.

\emph{Connected matroids with a series pair.}
Since \(M\) has no coloops, failure of cosimplicity means that
\(M\) has a series pair.
Choose \(N,f\) as in \autoref{lem:series}.
In case \(\mathrm{(i)}\) of \autoref{lem:series}, \(M\) is a
one-element series extension of \(N\) at \(f\).
We first bound \(b(N\backslash f)\) by counting circuits containing \(f\).
For each circuit \(C\) containing \(f\), fix a basis \(B_C\)
of \(N\backslash f\) extending \(C-f\).
Since \(f\) is not a coloop, \(B_C\) is also a basis of \(N\),
and \(C=C(f,B_C)\). Thus, \(C\mapsto B_C\) is injective.
The injection and \eqref{eq:oxley} applied to \(N\backslash f\) imply
\[
 b(N\backslash f)\ge d_f(N),\qquad
 b(N\backslash f)\ge2d(N\backslash f).
\]
Adding twice the first bound to the second, we obtain
\[
 3b(N\backslash f)\ge2d_f(N)+2d(N\backslash f)=2d(N).
\]
By the induction hypothesis and \autoref{lem:series}\(\mathrm{(i)}\),
\[
 2b(M)=2b(N)+2b(N\backslash f)
 \ge\Bigl(r(N)+\frac73\Bigr)d(N)
 >(r(M)+1)d(M).
\]

In case \(\mathrm{(ii)}\) of \autoref{lem:series},
\(M=P(N,M|\{f,x,y\})\). Suppose first that
\(d(N\backslash f)\ge1\).
By the two identities in \autoref{lem:series}\(\mathrm{(ii)}\),
\autoref{lem:pointed}, and \(d_f(N)=d(N)-d(N\backslash f)\), we deduce
\begin{align*}
 2b(M)-(r(M)+1)d(M)
 &=4b(N)+2b_f(N)-(r(N)+2)(d(N)+d_f(N)+1)\\
 &\ge2\bigl(2b(N)-(r(N)+1)d(N)\bigr)
       +r(N)\bigl(d(N\backslash f)-1\bigr).
\end{align*}
By induction, \(2b(N)-(r(N)+1)d(N)\ge0\), with equality only if
\(N\cong F_7\). Since \(d(N\backslash f)\ge1\), we have
\(r(N)(d(N\backslash f)-1)\ge0\). If \(N\cong F_7\), then \eqref{eq:fano} gives
\(r(N)(d(N\backslash f)-1)=3(7-1)=18>0\).
Hence, \(2b(M)>(r(M)+1)d(M)\) when \(d(N\backslash f)\ge1\).
If \(d(N\backslash f)=0\), then \(E(N)\) is a circuit by
\eqref{eq:circuit-case}. Hence, by \autoref{lem:series}\(\mathrm{(ii)}\),
\(b(M)=3r(N)+2\) and \(d(M)=3\), so
\[
 2b(M)-(r(M)+1)d(M)=3r(N)-2>0.
\]
Thus, the bound is strict in every case with \(r(M)\ge4\),
completing the induction.
\end{proof}

To prove \autoref{thm:deletion} in the cosimple case, we combine
\autoref{thm:rank} with the following bound comparing bases containing
and avoiding an element.

\begin{lemma}\label{lem:exchange-count}
Let \(M\) be a simple cosimple binary matroid with \(r(M)\ge3\).
For every \(e\in E(M)\),
\[
 2b_e(M)\le(r(M)-1)b(M\backslash e).
\]
\end{lemma}
\begin{proof}
We count basis exchanges that remove \(e\). By \eqref{eq:exchange}
and cosimplicity, each basis \(B\) containing \(e\) admits
\(|C^*(e,B)|-1\ge2\) such exchanges. In the reverse direction,
a basis \(B\) avoiding \(e\) admits \(|C(e,B)|-1\) exchanges that
insert \(e\), at most \(r(M)-1\) unless \(C(e,B)\) is spanning.
We compensate for each spanning-circuit exception by an extra exchange
from a distinct basis containing \(e\).
Let \(q\) be the number of spanning circuits containing \(e\).
For each such circuit \(C\), at least one of the bases
\(C-x\), with \(x\in C-e\), has \(|C^*(e,C-x)|>3\).
To see this, put \(B=C-e\) and suppose all these cocircuits
are triads. Since \(C^*(e,C-x)\) avoids \(C-\{e,x\}\),
\eqref{eq:binary-orthogonality} forces
\(C^*(e,C-x)\cap C=\{e,x\}\).
Thus, it meets \(B\) only in \(x\), so it is \(C^*(x,B)\).
Consequently,
\[
 C^*(e,C-x)=C^*(x,B)=\{e,x,z_x\},\qquad z_x\notin C.
\]
If \(z_x=z_y\) for distinct \(x,y\), the dual of
\eqref{eq:binary-symmetric-difference} implies that
\(\{x,y\}\) contains a cocircuit, contrary to cosimplicity.
Thus, the \(z_x\) are distinct. Each \(z_x\) belongs to
\(C^*(x,B)\) and to no other fundamental cocircuit of \(B\).
By \eqref{eq:exchange}, \(C(z_x,B)=\{z_x,x\}\), contrary
to simplicity.

Choose one such basis for each of the \(q\) spanning circuits.
These bases are distinct by \eqref{eq:spanning-bases}, so the
bases containing \(e\) supply at least \(q\) exchanges in addition
to their two exchanges guaranteed by cosimplicity.
Since \(e\) is not a coloop, the bases avoiding \(e\) are exactly
the bases of \(M\backslash e\).
Counting the exchanges in both directions by \eqref{eq:exchange}, we obtain
\begin{align*}
 2b_e(M)+q
 &\le\sum_{B\in\BB(M),\,e\in B} (|C^*(e,B)|-1)\\
 &=\sum_{B\in\BB(M),\,e\notin B} (|C(e,B)|-1)\\
 &\le(r(M)-1)b(M\backslash e)+q.
\end{align*}
In the last sum, precisely the \(q\) bases \(C-e\) have
\(|C(e,B)|=r(M)+1\); all other terms are at most \(r(M)-1\).
Subtracting \(q\) proves the assertion.
\end{proof}

\begin{proof}[Proof of \autoref{thm:deletion}]
We use strong induction on \(|E(M)|\), including the equality classification.
If \(E(M)\) is a circuit, then \(b(M\backslash e)=d(M)=1\).
Every rank-two matroid under consideration is \(U_{2,3}\),
so we may assume that \(E(M)\) is not a circuit and \(r(M)\ge3\).

\emph{Disconnected matroids.}
Write \(M=M_1\oplus M_2\), where both summands are nonempty
and \(e\in E(M_1)\).
Neither summand has coloops, so \(d(M_1),d(M_2)\ge1\).
By the induction hypothesis and \eqref{eq:oxley},
\[
 b(M\backslash e)=b(M_1\backslash e)b(M_2)\ge d(M_1)b(M_2)\ge2d(M_1)d(M_2)\ge d(M_1)+d(M_2)=d(M).
\]
Equality in the last inequality forces \(d(M_1)=d(M_2)=1\),
and the absence of coloops then makes \(E(M_2)\) its unique circuit.
Thus, \(b(M_2)=|E(M_2)|\ge3>2d(M_2)\), so the bound is strict.

\emph{Connected cosimple matroids.}
For \(M\not\cong F_7\), \autoref{thm:rank} and
\autoref{lem:exchange-count} imply
\[
 d(M)<\frac{2b(M)}{r(M)+1}
 =\frac{2(b_e(M)+b(M\backslash e))}{r(M)+1}
 \le b(M\backslash e).
\]
For \(F_7\), the deletion bound is strict by \eqref{eq:fano}.

\emph{Connected matroids with a series pair.}
Choose \(N,f\) as in \autoref{lem:series}.
In case \(\mathrm{(i)}\) of \autoref{lem:series}, \(M\) is a
one-element series extension of \(N\) at \(f\), replacing
\(f\) by \(\{x,y\}\).
Counting the bases that avoid \(e\) using the description following
\autoref{lem:series}\(\mathrm{(i)}\), we obtain
\begin{equation}\label{eq:series-deletions}
 b(M\backslash e)=
 \begin{cases}
 b(N\backslash f),&e\in\{x,y\},\\
 2b(N\backslash e)-b_f(N\backslash e),&e\in E(N)-f.
 \end{cases}
\end{equation}
By the induction hypothesis, the first value in
\eqref{eq:series-deletions} is at least \(d(N)=d(M)\);
the second is at least \(b(N\backslash e)\ge d(N)\), since
\(b_f(N\backslash e)\le b(N\backslash e)\).
Equality holds for \(e\in\{x,y\}\) precisely when
\(b(N\backslash f)=d(N)\). For \(e\in E(N)-f\), equality holds precisely when
\[
 b(N\backslash e)=d(N)
 \quad\text{and}\quad f\text{ is a coloop of }N\backslash e.
\]
Here \(b_f(N\backslash e)=b(N\backslash e)\) holds exactly when
\(f\) is a coloop of \(N\backslash e\).
Since \(N\) has no coloops, the coloop condition is equivalent
to \(\{e,f\}\) being a series pair of \(N\), by the cocircuit
description of deletion \cite[(3.1.16)]{Oxley2011}.

To determine equality in \eqref{eq:series-deletions}, first note that
\(E(N)\) cannot be a circuit, since then \(E(M)\) would
also be a circuit. By the induction hypothesis, equality
requires \(N=P(N|C_1,N|C_2)\), where \(C_1\) is a circuit
and \(C_2\) is a triangle.
For \(e\in\{x,y\}\), equality requires \(f\notin C_2\).
For \(e\in E(N)-f\), equality is possible only if \(e\notin C_2\)
and \(\{e,f\}\) is a series pair. By \eqref{eq:binary-orthogonality} with \(C_2\),
we also have \(f\notin C_2\).
Thus, \eqref{eq:series-circuits} shows that the extension replaces \(C_1\) by
\((C_1-f)\cup\{x,y\}\), preserves \(C_2\), and leaves \(e\)
outside \(C_2\).

In case \(\mathrm{(ii)}\) of \autoref{lem:series},
\(M=P(N,M|\{f,x,y\})\). By the circuit count in that lemma,
\begin{equation}\label{eq:triangle-circuits}
 d(M)=2d(N)-d(N\backslash f)+1.
\end{equation}
From the basis description following \autoref{lem:series}\(\mathrm{(ii)}\),
we compute
\begin{equation}\label{eq:triangle-deletions}
 b(M\backslash e)=
 \begin{cases}
 b(N),&e\in\{x,y\},\\
 b(N)+b(N\backslash f),&e=f,\\
 2b(N\backslash e)+b_f(N\backslash e),&e\in E(N)-f.
 \end{cases}
\end{equation}
If \(d(N\backslash f)\ge1\), then \eqref{eq:oxley} and
\eqref{eq:triangle-circuits} imply
\[
 b(N)\ge2d(N)\ge d(M).
\]
Equality would require \(N\cong F_7\) and
\(d(N\backslash f)=1\), contrary to \eqref{eq:fano}.
Thus, \(b(M\backslash e)>d(M)\) for \(e\in\{x,y,f\}\)
by \eqref{eq:triangle-deletions}.
For \(e\in E(N)-f\), the element \(f\) is a nonloop of
\(N\backslash e\), so \(b_f(N\backslash e)\ge1\).
The induction hypothesis and \eqref{eq:triangle-deletions}
imply \(b(M\backslash e)\ge2d(N)+1>d(M)\).

If \(d(N\backslash f)=0\), \eqref{eq:circuit-case} makes
\(E(N)\) a circuit. Then \(d(M)=3\), and \(b(M\backslash e)\) equals
\(|E(N)|\), \(|E(N)|+1\), or \(3\), according to the three cases in
\eqref{eq:triangle-deletions}.
Thus, equality holds for every \(e\in E(N)-f\),
and also for \(e\in\{x,y\}\) when \(N\) is a triangle.
When \(N\) is a triangle and \(e\in\{x,y\}\), interchange the two triangles.

Conversely, equality holds when \(E(M)\) is a circuit. A parallel connection
of a circuit \(C_1\) and a triangle \(C_2\) has three circuits
by \eqref{eq:parallel-circuits}.
Taking \(N=M|C_1\) in \eqref{eq:triangle-deletions}, we count
exactly three bases after deleting any element outside \(C_2\).
This completes the equality classification and the induction.
\end{proof}

\begin{proof}[Proof of \autoref{cor:growth}]
The case \(s=0\) is immediate, so assume \(s\ge1\).
Take successive matroids \(N_0,\ldots,N_s\), with \(N_0=N\) and \(N_s=M\),
where \(N_i\) is obtained from \(N_{i-1}\) by a one-element series
extension at \(f_i\) for \(1\le i\le s\).
Every \(N_i\) is simple and binary and has no coloops.
By the series-extension identities in \autoref{lem:series}\(\mathrm{(i)}\)
and \autoref{thm:deletion},
\begin{equation}\label{eq:growth-step}
 b(N_i)-b(N_{i-1})=b(N_{i-1}\backslash f_i)
 \ge d(N_{i-1})=d(N).
\end{equation}
Summing proves the bound. If \(N\) is cosimple,
the inequality in \eqref{eq:growth-step} is strict for \(i=1\).

Suppose equality holds. Then every step in \eqref{eq:growth-step}
is an equality. In particular, \autoref{thm:deletion} applied to
\(N_{s-1}\) and \(f_s\) shows that \(E(N_{s-1})\) is a circuit, or that
\(N_{s-1}\) is a parallel connection of a circuit and a triangle with
\(f_s\) outside the triangle. The series extension \(M\) has the same
form by \eqref{eq:series-circuits}.

Conversely, if \(E(M)\) is a circuit, \eqref{eq:series-circuits} shows that
every \(E(N_i)\) is a circuit. Hence,
\(b(N_{i-1}\backslash f_i)=d(N)=1\) at every step.
For the parallel-connection case, we trace the extensions backwards
to verify equality at every step.
Suppose \(M=P(M|C_1,M|C_2)\), where \(C_1\) is a circuit
and \(C_2\) is a triangle meeting \(C_1\) at \(f\).
Let \(\{x,y\}\) be the series pair introduced in the last extension.
Up to relabelling \(y\), its predecessor \(N_{s-1}\) is \(M/x\).
By \eqref{eq:binary-orthogonality}, the pair lies in
\(C_1-f\) or \(C_2-f\). The latter is impossible, since
\(C_2-x\) would be a two-element circuit of the simple matroid \(M/x\).
Thus, \(\{x,y\}\subseteq C_1-f\). By
\cite[Proposition~7.1.15(v)]{Oxley2011},
\(M/x=P((M|C_1)/x,M|C_2)\).
The first factor has the unique circuit \(C_1-x\), which has at least
three elements because \(M/x\) is simple.
Hence, \(N_{s-1}\) has the same form, with \(C_2\) unchanged
and the relabelled element \(y\) outside \(C_2\).
Repeating the argument shows that every forward step extends an
element outside \(C_2\) and attains equality in \eqref{eq:growth-step}
by \autoref{thm:deletion}.
\end{proof}
\section*{Acknowledgements}
This work is supported by the Guangdong Basic and Applied Basic Research Foundation (Grant No. 2026A1515012237).
\section*{Declaration of AI use}
The author developed the results and original proofs.
The author used Codex (OpenAI) to assist in simplifying the original proofs, clarifying the exposition, and polishing the language.
The author takes full responsibility for the content of the manuscript.

\begin{thebibliography}{9}
\setlength{\itemsep}{2pt}
\setlength{\parsep}{0pt}
\setlength{\parskip}{0pt}
\bibitem{Cho2026}
C. H. Cho,
\emph{Circuits and enumeration problems for matroids},
Ph.D. dissertation, Louisiana State University, Baton Rouge, LA, 2026.
\bibitem{EllisLeader2018}
D. Ellis and I. Leader,
An isoperimetric inequality for antipodal subsets of the discrete cube,
\emph{European J. Combin.} \textbf{70} (2018), 149--154.
\bibitem{Oxley1983}
J. G. Oxley,
On the numbers of bases and circuits in simple binary matroids,
\emph{European J. Combin.} \textbf{4} (1983), 169--178.
\bibitem{Oxley2011}
J. G. Oxley,
\emph{Matroid Theory}, second edition,
Oxford Graduate Texts in Mathematics 21,
Oxford University Press, Oxford, 2011.
\bibitem{OxleyWang2019}
J. G. Oxley and S. Wang,
Dependencies among dependencies in matroids,
\emph{Electron. J. Combin.} \textbf{26} (2019), no.~3, Paper No.~P3.46, 12~pp.
\bibitem{Welsh1976}
D. J. A. Welsh,
\emph{Matroid Theory},
London Mathematical Society Monographs 8,
Academic Press, London, 1976.
\end{thebibliography}

\end{document}
